\subsection{The Exterior Algebra}
\label{linalg:sec:the_exterior_algebra}

This section is meant to give a quick overview of the definition and most essential properties of the exterior algebra $\Lambda^\ast V$ of a finite dimensional vector space $V$ (over some field $\mathbb{K}$). To define the $k$-fold exterior power $\Lambda^k V$ consider the permutation maps 
\[
P_\sigma \colon V^{\otimes k} \to V^{\otimes k} 
\]
for $\sigma \in S(k) \footnote{S(k) \text{ denotes the symmetric group on $n$ elements}}$ given by 
\[
P_\sigma(v_1\otimes \ldots \otimes v_k) := \mathrm{sgn}(\sigma)v_{\sigma(1)}\otimes \ldots \otimes v_{\sigma(k)}.
\]
These are well-defined by the universal property of the tensor product. We then call a vector $\alpha \in V^{\otimes k}$ \textbf{alternating} if 
\[
P_\sigma(\alpha) = \alpha \quad \forall \sigma \in S(k)
\]
and define $\Lambda^k V \subseteq V^{\otimes k}$ as the subspace of all alternating vectors. Given any vector in $V^{\otimes k}$ there is a canonical way to make it into an alternating one using the \textbf{antisymmetrization map}
\[
\mathrm{Alt} \colon V^{\otimes k} \to \Lambda^k V, \quad \alpha \mapsto \frac{1}{k!}\sum_{\sigma \in S(k)} P_\sigma(\alpha).
\]
Note that this map is surjective and $\mathrm{Alt}(\alpha) = \alpha$ if and only if $\alpha$ is alternating. Using the antisymmetrization map we can define the \textbf{exterior product} (or \textbf{wedge product}) as the bilinear map 
\begin{align*}
\Lambda^k V \times \Lambda^l V & \to \Lambda^{k+l} V, \\
(\alpha,\beta) &\mapsto \alpha \wedge \beta := \frac{(k+l)!}{k!l!}\mathrm{Alt}(\alpha \otimes \beta).
\end{align*}
Without proof we state the facts, that this product is associative and graded commutative. Moreover if $e_1,\ldots,e_n \in V$ is a basis of $V$ then 
\[
\{e_{i_1} \wedge \ldots \wedge e_{i_k}\}_{1 \leq i_1 \leq \ldots \leq i_k \leq n}
\]
is a basis of $\Lambda^k V$. If we identify $\Lambda^k V$ with the space of alternating maps $V^\ast \times \ldots \times V^\ast \to \R$, these basis vectors satisfy 
\[
(e_{i_1} \wedge \ldots \wedge e_{i_k})(\epsilon^{j_1},\ldots,\epsilon^{j_k}) = \begin{dcases}
    1 &, (i_1,\ldots,i_k) = (j_1,\ldots,j_k) \\
    0 & \text{otherwise}
\end{dcases}
\]
where $\epsilon^1, \ldots, \epsilon^n \in V^\ast$ is the dual basis of $e_1,\ldots,e_n$ and $1 \leq j_1 < \ldots < j_k \leq n$ is any ordered tuple. This is the reason for the normalization factor in the definition of the wedge product. 
The exterior algebra is then defined as 
\[
\Lambda^\ast V := \bigoplus_{k=0}^\infty \Lambda^k V
\]
with multiplication given by the wedge product (we set $V^{\otimes 0} = \Lambda^0 V := \mathbb{K}$).

We can also characterize the exterior power and the exterior algebra by means of a universal property. 

\begin{definition}[Universal property of the exterior power]
    \label{linalg:def:universal_property_exterior_power}
    The ($k$-fold) exterior power of $V$, denoted $\Lambda^k V$, is a vector space together with an alternating multilinear map 
    \[
    \pi \colon V^{\times k} \to \Lambda^k V,
    \]
    such that for any other alternating multilinear map 
    \[
    T \colon V^{\times k} \to W
    \]
    into a vector space $W$ there exists a unique linear map 
    \[
    \widetilde{T} \colon \Lambda^k V \to W
    \]
    making the following diagram commute 
    \[\begin{tikzcd}[column sep = 3em]
	{V^{\times k}} & W \\
	{\Lambda^k V}
	\arrow["T", from=1-1, to=1-2]
	\arrow["\pi"', from=1-1, to=2-1]
	\arrow["\widetilde{T}"', dashed, end anchor = {[xshift = 5, yshift = -3]}, from=2-1, to=1-2]
\end{tikzcd}\]
\end{definition}

\begin{definition}[Universal property of the exterior algebra]
    \label{linalg:def:universal_property_exterior_algebra}
    The exterior algebra of $V$, denoted $\Lambda^\ast V$, is an algebra together with a linear map 
    \[
    \iota \colon V \to \Lambda^\ast V,
    \]
    which satisfies $\iota(v)^2 = 0$ for all $v\in V$ and such that for any linear map 
    \[
    \psi \colon V \to A^\prime
    \]
    into an algebra $A^\prime$ with $\psi(v)^2 = 0$ for all $v\in V$ there exists a unique algebra homomorphism
    \[
    L \colon \Lambda^\ast V \to A^\prime
    \]
    making the following diagram commute 
    \[\begin{tikzcd}[column sep = 3em]
	{V} & \Lambda^\ast V\\
	& {A^\prime}
	\arrow["\iota", from=1-1, to=1-2]
	\arrow["\psi"', from=1-1, to=2-2]
	\arrow["L", dashed, from=1-2, to=2-2]
\end{tikzcd}\]
\end{definition}


\begin{remark}
    The universal property of the exterior power is analogous to the one of the tensor product, where we think of $\pi$ as the product map $(v_1, \ldots, v_k) \mapsto v_1 \wedge \ldots \wedge v_k$. So alternating multilinear maps on $V^{\times k}$ are in $1:1$ correspondence to linear maps on $\Lambda^k V$.
    In Definition \ref{linalg:def:universal_property_exterior_algebra} we think of $\iota$ as an inlcusion and of $L$ as extending the map $\psi$.
\end{remark}

\begin{lemma}
    \label{linalg:lem:universal_property_exterior_power}
    The construction of $\Lambda^k V \subseteq V^{\otimes k}$ at the beginning of the section together with the product map $\pi \colon V^{\times k} \to \Lambda^k V, \; (v_1, \ldots, v_k) \mapsto v_1 \wedge \ldots \wedge v_k$ satisifes the universal property in Definition \ref{linalg:def:universal_property_exterior_power}.
\end{lemma}

\begin{proof}
    First of all it follows from the graded commutativity of the wedge product, that $\pi$ is alternating. Uniqueness of $\widetilde{T}$ is also clear. To show existence choose a basis $e_1,\ldots,e_n \in V$ of $V$. As mentioned above this yields a basis of $\Lambda^kV$ and we define $\widetilde{T}$ on this basis (necessarily) by 
    \[
    \widetilde{T}(e_{i_1} \wedge \ldots \wedge e_{i_k}) := T(e_{i_1},  \ldots, e_{i_k})
    \]
    for $1 \leq i_1 < \ldots < i_k \leq n$. To see that 
    \[\begin{tikzcd}[column sep = 3em]
	{V^{\times k}} & W \\
	{\Lambda^k V}
	\arrow["T", from=1-1, to=1-2]
	\arrow["\pi"', from=1-1, to=2-1]
	\arrow["\widetilde{T}"', dashed, end anchor = {[xshift = 5, yshift = -3]}, from=2-1, to=1-2]
\end{tikzcd}\]
does in fact commute it suffices to verify commutativity on basis vectors $(e_{i_1}, \ldots, e_{i_k})$ of $V^{\times k}$ (where the tuple $(i_1,\ldots,i_k)$ is not necessarily ordered!). But since $T$ and $\pi$ are both alternating we may assume after a suitable permutation of $(i_1,\ldots,i_k)$, that this tuple is ordered, such that the statement follows from the definiton of $\widetilde{T}$.
\end{proof}

\begin{lemma}
    The construction of $\Lambda^\ast V$ at the beginning of the section together with the canonical inclusion map $V \hookrightarrow \Lambda^\ast V$ satisifes the universal property in Definition \ref{linalg:def:universal_property_exterior_algebra}.
\end{lemma}

\begin{proof}
    Again uniqueness of $L$ is clear. For the existence of $L$ we define it on $\Lambda^k V$ (necessarily) by 
    \[
    L(v_1 \wedge \ldots \wedge v_k) := \prod_{i=1}^k \psi^\prime(v_i)
    \]
    using the universal property of the exterior power shown in Lemma \ref{linalg:lem:universal_property_exterior_power}.
\end{proof}


To conclude this section let us mention another commonly used construction of the exterior power and exterior algebra. To this end let 
\[
S^k(V) := \mathrm{span}\{v_1\otimes\ldots\otimes v_k \mid \exists \; 1 \leq i,j \leq k, i \neq j \text{ with } v_i = v_j\} \subseteq V^{\otimes k}
\]
Then it is not hard to show that $\modulo{V^{\otimes k}}{S^k(V)}$ together with the canonical map 
\[
V^{\times k} \to V^{\otimes k} \to \modulo{V^{\otimes k}}{S^k(V)}
\]
also satisifies the universal property in Definition \ref{linalg:def:universal_property_exterior_power}. Hence this is also a valid and commonly used construction of the exterior power. With this definition 
\[
\alpha \wedge \beta = [\alpha \otimes \beta].
\]


% In particular this implies that there exists a unique isomorphism 
% \[
% \Lambda^k V \xrightarrow{\cong} \modulo{V^{\otimes k}}{S^k(V)}
% \]
% making the following diagram commute 
% \[\begin{tikzcd}[column sep = 3em]
% 	{V^{\otimes k}} & \modulo{V^{\otimes k}}{S^k(V)} \\
% 	{\Lambda^k V}
% 	\arrow[from=1-1, to=1-2]
% 	\arrow["\mathrm{Alt}"', from=1-1, to=2-1]
% 	\arrow["\cong", end anchor = {[xshift = 5, yshift = -3]}, from=2-1, to=1-2]
% \end{tikzcd}\]
% (where $\Lambda^k V \subseteq V^{\otimes k}$ denotes our original construction of the exterior power given at the beginning of the section). This gives us an explicit description of the kernel of $\mathrm{Alt}$, namely 
% \[
% \mathrm{ker(Alt)} = S^k(V).
% \]


\noindent Similiarly we can also define the exterior algebra as a quotient of the tensor algebra
\[
T(V) := \bigoplus_{k=0}^\infty V^{\otimes k},
\]
by dividing out the two sided ideal $I \subseteq T(V)$ generated by elements of the form $v \otimes v \in T(V)$ for $v \in V$, i.e. 
\[
\Lambda^\ast V := \modulo{T(V)}{I}.
\]
Together with the canonical map $V \to \Lambda^\ast V$ this definition is also easily seen to satisfy the universal property in Definition \ref{linalg:def:universal_property_exterior_algebra}.
